\chapter{作业图片转写待审核稿}

\begin{dxtips}

本文件为图片与 Markdown 笔记整理出的待审核转写稿，已尽量按 elegantbook 环境归类；正式并入正文前仍建议逐条复核。

\end{dxtips}

\section*{1}

\begin{note}[来源上级主题]

3.8作业 (3\%208\%E4\%BD\%9C\%E4\%B8\%9A\%20206d2efa3f42807aaea0cdba1c8e0355.md)

\end{note}

\begin{homework}

4r 里面 8 16 24

(4)包括加法 所以 4 8 12

4生成理想是通过主环偶数环生成的所以4r=（4）

例三2，x生成的环就是 2p1(x)+xp2(x) 主理想得是一个元生成的所以设是px

让x=ah(x)

2=ap(x)那么p(x)=1/a就可以了 hx=1/ax也是可以的

\end{homework}

\begin{note}[关联图片]
以下图片已纳入本条整理，当前版本优先依据 Markdown 文字整理，图片细节可在审核时对照：

1. image.png\\
2. image 1.png\\
3. image 2.png\\
4. image 3.png\\
5. image 4.png\\

\end{note}

\begin{note}[来源路径]

近世代数/近世代数/无标题/1 206d2efa3f42800a84cefbb4eef7fd0b.md

\end{note}

\section*{1欧式环}

\begin{note}[来源上级主题]

第四章作业 (\%E7\%AC\%AC\%E5\%9B\%9B\%E7\%AB\%A0\%E4\%BD\%9C\%E4\%B8\%9A\%20204d2efa3f4280b4b282f638dd2dce07.md)

\end{note}

\begin{homework}

证明域都是 欧式环域内非零元都是有逆元所以q=a/b恒成立

a=qb+r r=0恒成立所以可以

把域中所有元素都映射到1

欧式环定义

一个整环里面首先构建一个非零元到非负整数的映射

不等于0的元a 和任意元b都能写成

b=qa+r r=0 或者ϕ(r)<ϕ(a) 这里就是让ba的逆写成q可以保证r=0恒成立还有因为是域每一个数都有逆元

\end{homework}

\begin{note}[关联图片]
以下图片已纳入本条整理，当前版本优先依据 Markdown 文字整理，图片细节可在审核时对照：

1. image.png\\

\end{note}

\begin{note}[来源路径]

近世代数/近世代数/无标题/1欧式环 204d2efa3f42806fb730dc64fc6e2b23.md

\end{note}

\section*{2}

\begin{note}[来源上级主题]

3.8作业 (3\%208\%E4\%BD\%9C\%E4\%B8\%9A\%20206d2efa3f42807aaea0cdba1c8e0355.md)

\end{note}

\begin{homework}

只要理想组合起来等于主理想即可

\end{homework}

\begin{note}[图片 OCR 摘录，待人工复核]
以下内容来自源图片识别，便于审核时对照：

\textbf{图片 1}（image.png）
\begin{Verbatim}[fontsize=\small]
3. BER IM REAMAIM. TEM, At (2, x)
fE—P ERI,  |

fi HREAMBR, WAR x ISHAM S, ATT
CMM O, =) MHL 2=1, Al (2, ») SF aw
fA C1).
\end{Verbatim}

\end{note}

\begin{note}[来源路径]

近世代数/近世代数/无标题/2 20ad2efa3f4280f18cc9e1642c615bfd.md

\end{note}

\section*{2正规子群}

\begin{note}[来源上级主题]

第二章作业 (\%E7\%AC\%AC\%E4\%BA\%8C\%E7\%AB\%A0\%E4\%BD\%9C\%E4\%B8\%9A\%2020bd2efa3f42805f9ecef7b781a02d72.md)

\end{note}

\begin{homework}

2. WH, WRRETFRNRERCERETH.

fe AN AN. BAG HRTAREFRBANINN:,

EGHW—-A+*FH(88. YBl2). RN HIER, NiNN,.-

FEG IW —- ARETE. SaG, nC N,NN., MAnEN,,

nEN,, (AN, MN, RAE STH, PUM ana'eNn,,

‘ana'EN,, Ait

ana'E€N, Nn,

FREHEH 2, NNN. —AHETH.

\end{homework}

\begin{note}[图片 OCR 摘录，待人工复核]
以下内容来自源图片识别，便于审核时对照：

\textbf{图片 1}（image.png）
\begin{Verbatim}[fontsize=\small]
2. WH, WRRETFRNRERCERETH.

fe AN AN. BAG HRTAREFRBANINN:,
EGHW—-A+*FH(88. YBl2). RN HIER, NiNN,.-
FEG IW —- ARETE. SaG, nC N,NN., MAnEN,,
nEN,, (AN, MN, RAE STH, PUM ana'eNn,,
‘ana'EN,, Ait

ana'E€N, Nn,

FREHEH 2, NNN. —AHETH.
\end{Verbatim}

\end{note}

\begin{note}[来源路径]

近世代数/近世代数/无标题/2正规子群 20bd2efa3f4280879ad4e4e896b6e022.md

\end{note}

\section*{3.2整环}

\begin{note}[来源上级主题]

第三章作业 (\%E7\%AC\%AC\%E4\%B8\%89\%E7\%AB\%A0\%E4\%BD\%9C\%E4\%B8\%9A\%20205d2efa3f428004b442cf8cf50564e6.md)

\end{note}

\begin{homework}

子级 项目: 第一题 (\%E7\%AC\%AC\%E4\%B8\%80\%E9\%A2\%98\%20209d2efa3f4280518e01dd4451714875.md), 无标题 (\%E6\%97\%A0\%E6\%A0\%87\%E9\%A2\%98\%20209d2efa3f4280769b8efec54d64c727.md)

\end{homework}

\begin{note}[来源路径]

近世代数/近世代数/无标题/3 2整环 209d2efa3f4280cda320c843bba252f6.md

\end{note}

\section*{3.3}

\begin{note}[来源上级主题]

第三章作业 (\%E7\%AC\%AC\%E4\%B8\%89\%E7\%AB\%A0\%E4\%BD\%9C\%E4\%B8\%9A\%20205d2efa3f428004b442cf8cf50564e6.md)

\end{note}

\begin{homework}

子级 项目: 第三题 (\%E7\%AC\%AC\%E4\%B8\%89\%E9\%A2\%98\%20209d2efa3f4280918485e1eeb5bf0217.md), 无零因子作用 (\%E6\%97\%A0\%E9\%9B\%B6\%E5\%9B\%A0\%E5\%AD\%90\%E4\%BD\%9C\%E7\%94\%A8\%2020ad2efa3f428015bfb2db4f02684706.md)

\end{homework}

\begin{note}[关联图片]
以下图片已纳入本条整理，当前版本优先依据 Markdown 文字整理，图片细节可在审核时对照：

1. image.png\\
2. image 1.png\\

\end{note}

\begin{note}[来源路径]

近世代数/近世代数/无标题/3 3 209d2efa3f428020b15ef2f886bfb11e.md

\end{note}

\section*{3.4}

\begin{note}[来源上级主题]

第三章作业 (\%E7\%AC\%AC\%E4\%B8\%89\%E7\%AB\%A0\%E4\%BD\%9C\%E4\%B8\%9A\%20205d2efa3f428004b442cf8cf50564e6.md)

\end{note}

\begin{homework}

子级 项目: 4元群 (4\%E5\%85\%83\%E7\%BE\%A4\%20209d2efa3f42801c9f8ef20ec54114ef.md)

\end{homework}

\begin{note}[来源路径]

近世代数/近世代数/无标题/3 4 209d2efa3f4280eca043fa61064594c9.md

\end{note}

\section*{3.5作业}

\begin{note}[来源上级主题]

第三章作业 (\%E7\%AC\%AC\%E4\%B8\%89\%E7\%AB\%A0\%E4\%BD\%9C\%E4\%B8\%9A\%20205d2efa3f428004b442cf8cf50564e6.md)

\end{note}

\begin{homework}

第一题逆元

(a+(-a))x=ax+-ax=0.  a*b=1 baxb=bxab  bx=xb 交换ax可以让a与b挨着

-ax=-xa                   证明如果还有一个域就是本身 可以从生成元考虑，任意a+bi都能组合出来

\end{homework}

\begin{note}[关联图片]
以下图片已纳入本条整理，当前版本优先依据 Markdown 文字整理，图片细节可在审核时对照：

1. image.png\\
2. image 1.png\\

\end{note}

\begin{note}[来源路径]

近世代数/近世代数/无标题/3 5作业 209d2efa3f4280eb969cc9dce9865299.md

\end{note}

\section*{3.8作业}

\begin{note}[来源上级主题]

第三章作业 (\%E7\%AC\%AC\%E4\%B8\%89\%E7\%AB\%A0\%E4\%BD\%9C\%E4\%B8\%9A\%20205d2efa3f428004b442cf8cf50564e6.md)

\end{note}

\begin{homework}

子级 项目: 1 (1\%20206d2efa3f42800a84cefbb4eef7fd0b.md), 2 (2\%2020ad2efa3f4280f18cc9e1642c615bfd.md)

\end{homework}

\begin{note}[来源路径]

近世代数/近世代数/无标题/3 8作业 206d2efa3f42807aaea0cdba1c8e0355.md

\end{note}

\section*{3.9作业}

\begin{note}[来源上级主题]

第三章作业 (\%E7\%AC\%AC\%E4\%B8\%89\%E7\%AB\%A0\%E4\%BD\%9C\%E4\%B8\%9A\%20205d2efa3f428004b442cf8cf50564e6.md)

\end{note}

\begin{homework}

子级 项目: 第一题 (\%E7\%AC\%AC\%E4\%B8\%80\%E9\%A2\%98\%20205d2efa3f4280d38576e58476d8c41d.md), 第二题 (\%E7\%AC\%AC\%E4\%BA\%8C\%E9\%A2\%98\%20205d2efa3f4280d18289e2c61ebcba54.md), 第三题再看 (\%E7\%AC\%AC\%E4\%B8\%89\%E9\%A2\%98\%E5\%86\%8D\%E7\%9C\%8B\%20205d2efa3f42808f87efe6f9cc3fa0d4.md), 第三题 (\%E7\%AC\%AC\%E4\%B8\%89\%E9\%A2\%98\%20205d2efa3f4280c295f6ed5b008ce795.md)

\end{homework}

\begin{note}[来源路径]

近世代数/近世代数/无标题/3 9作业 205d2efa3f428077b8a0d0acdba61717.md

\end{note}

\section*{3，4}

\begin{note}[来源上级主题]

第三周作业 (\%E7\%AC\%AC\%E4\%B8\%89\%E5\%91\%A8\%E4\%BD\%9C\%E4\%B8\%9A\%2020bd2efa3f428003b74ac10b654adc8a.md)

\end{note}

\begin{homework}

3. RECHR-TWEBANARH. EGRET 2

ER) 7G EK) YX — 5 BF BE

RB HTB 2A, GEMAT 2 NHRHTREBR. (8

GRA-*WHIKH 7, RRB Ie. TRATCHME

BR, BCBMST 2H THREAD.

\end{homework}

\begin{note}[图片 OCR 摘录，待人工复核]
以下内容来自源图片识别，便于审核时对照：

\textbf{图片 1}（image.png）
\begin{Verbatim}[fontsize=\small]
3. RECHR-TWEBANARH. EGRET 2
ER) 7G EK) YX — 5 BF BE

RB HTB 2A, GEMAT 2 NHRHTREBR. (8
GRA-*WHIKH 7, RRB Ie. TRATCHME
BR, BCBMST 2H THREAD.
\end{Verbatim}

\textbf{图片 2}（image 1.png）
\begin{Verbatim}[fontsize=\small]
4. — ARIS — Toc a PR
8 SCGR-TARRM CRGHE-UR, BA
17
\end{Verbatim}

\textbf{图片 3}（image 2.png）
\begin{Verbatim}[fontsize=\small]
a, a®, a’, >.   fa
AREAS. UGTA 7, i,.f>i, ai=al. FA
ai FPG, 48         i-
(1)         aiize
OPE, FEETEWAM 7, OE (1) RO. ALE TE RN
\end{Verbatim}

\textbf{图片 4}（image 3.png）
\begin{Verbatim}[fontsize=\small]
Y oEG
gla. - oe ~ ae
nize
heekeg)-jfeie 1 e
Mme oa ep febell|
\end{Verbatim}

\end{note}

\begin{note}[来源路径]

近世代数/近世代数/无标题/3，4 20bd2efa3f4280daa104ec7e9deedc65.md

\end{note}

\section*{4}

\begin{note}[来源上级主题]

第五周作业 (\%E7\%AC\%AC\%E4\%BA\%94\%E5\%91\%A8\%E4\%BD\%9C\%E4\%B8\%9A\%2020bd2efa3f428007923bd76a797a9ede.md)

\end{note}

\begin{homework}

4. RECHREAH, FH\#EGHGHA. TAG thE

RR,

4  HFG5@GAA, Gie—T+ RR. HG =(a), io

EGAGNASWH OT, o—o, HEMERME, BA

OP, Ha"eG, fia"—>g, \{A a"—>a . Hfig=a.

\end{homework}

\begin{note}[图片 OCR 摘录，待人工复核]
以下内容来自源图片识别，便于审核时对照：

\textbf{图片 1}（image.png）
\begin{Verbatim}[fontsize=\small]
4. RECHREAH, FH#EGHGHA. TAG thE
RR,

4  HFG5@GAA, Gie—T+ RR. HG =(a), io
EGAGNASWH OT, o—o, HEMERME, BA
OP, Ha"eG, fia"—>g, {A a"—>a . Hfig=a.
\end{Verbatim}

\textbf{图片 2}（image 1.png）
\begin{Verbatim}[fontsize=\small]
XE, GHB—TBE 0 HPI G =).
\end{Verbatim}

\end{note}

\begin{note}[来源路径]

近世代数/近世代数/无标题/4 20bd2efa3f4280efba7cd41e90a2076b.md

\end{note}

\section*{循环群}

\begin{note}[来源上级主题]

第五周作业 (\%E7\%AC\%AC\%E4\%BA\%94\%E5\%91\%A8\%E4\%BD\%9C\%E4\%B8\%9A\%2020bd2efa3f428007923bd76a797a9ede.md)

\end{note}

\begin{homework}

由一个元形成群所有元素 所有元素都可以写成生成元的次方

1=nq+mp n,m为整数

呜呜呜呜呜

\end{homework}

\begin{note}[关联图片]
以下图片已纳入本条整理，当前版本优先依据 Markdown 文字整理，图片细节可在审核时对照：

1. image.png\\
2. image 1.png\\
3. image 2.png\\
4. image 3.png\\
5. image 4.png\\
6. image 5.png\\

\end{note}

\begin{note}[来源路径]

近世代数/近世代数/无标题/循环群 20bd2efa3f4280d99b3eea67ac6b65ad.md

\end{note}

\section*{循环群的商群还是循环群}

\begin{note}[来源上级主题]

第二章作业 (\%E7\%AC\%AC\%E4\%BA\%8C\%E7\%AB\%A0\%E4\%BD\%9C\%E4\%B8\%9A\%2020bd2efa3f42805f9ecef7b781a02d72.md)

\end{note}

\begin{homework}

\% G=(a) Hr N<G

TRA ZR, ERR TRA TE, Be

NG.

VbeG,b=a bN=aN=(aN)’

Bek «= G/N =\{bN |b G\} = (aN) i)

\end{homework}

\begin{note}[图片 OCR 摘录，待人工复核]
以下内容来自源图片识别，便于审核时对照：

\textbf{图片 1}（image.png）
\begin{Verbatim}[fontsize=\small]
% G=(a) Hr N<G
TRA ZR, ERR TRA TE, Be
NG.
VbeG,b=a bN=aN=(aN)’
Bek «= G/N ={bN |b G} = (aN) i)
\end{Verbatim}

\textbf{图片 2}（image 1.png）
\begin{Verbatim}[fontsize=\small]
RNB GI —TAEFHN. JEN WTA BRE
PRE
S={aN, bN, cN, +}
FMB, NY
(xN) CyN) =+ary) N
\end{Verbatim}

\end{note}

\begin{note}[来源路径]

近世代数/近世代数/无标题/循环群的商群还是循环群 20bd2efa3f4280278080f5000566e0e2.md

\end{note}

\section*{循环群的子群还是循环群}

\begin{note}[来源上级主题]

第二章作业 (\%E7\%AC\%AC\%E4\%BA\%8C\%E7\%AB\%A0\%E4\%BD\%9C\%E4\%B8\%9A\%2020bd2efa3f42805f9ecef7b781a02d72.md)

\end{note}

\begin{homework}

其实也就是要证明子群中最小的正幂次方ak是生成元，设其他的元素不整除则m=kq+r群有逆元 所以ar也在循环群内 因为r比k小，k又是最小的正幂r只能等于0所有元都是k的整数倍

\end{homework}

\begin{note}[关联图片]
以下图片已纳入本条整理，当前版本优先依据 Markdown 文字整理，图片细节可在审核时对照：

1. image.png\\

\end{note}

\begin{note}[来源路径]

近世代数/近世代数/无标题/循环群的子群还是循环群 20bd2efa3f42808aba03cfbc6170a459.md

\end{note}

\section*{找子群再看明天早上}

\begin{note}[来源上级主题]

第二章作业 (\%E7\%AC\%AC\%E4\%BA\%8C\%E7\%AB\%A0\%E4\%BD\%9C\%E4\%B8\%9A\%2020bd2efa3f42805f9ecef7b781a02d72.md)

\end{note}

\begin{homework}

1. RS AAT H.

RW S,BRAUETLR:

Ss\#4; (D)=\{(\}s

((12)) =\{(12), (1) \}s

((13)) = \{ (13), (1) \}s

((23)) = \{ (23), (1) \}s

((123)) = \{ (123), (132), (1)\},0

\end{homework}

\begin{note}[图片 OCR 摘录，待人工复核]
以下内容来自源图片识别，便于审核时对照：

\textbf{图片 1}（image.png）
\begin{Verbatim}[fontsize=\small]
1. RS AAT H.

RW S,BRAUETLR:
Ss#4; (D)={(}s
((12)) ={(12), (1) }s
((13)) = { (13), (1) }s
((23)) = { (23), (1) }s
((123)) = { (123), (132), (1)},0
\end{Verbatim}

\end{note}

\begin{note}[来源路径]

近世代数/近世代数/无标题/找子群再看明天早上 20bd2efa3f4280e39200d13ff510ae80.md

\end{note}

\section*{无标题}

\begin{note}[来源上级主题]

第四周作业 (\%E7\%AC\%AC\%E5\%9B\%9B\%E5\%91\%A8\%E4\%BD\%9C\%E4\%B8\%9A\%2020bd2efa3f4280a583d6e86ac317e9a1.md)

\end{note}

\begin{homework}

待根据图片进一步人工校对。

\end{homework}

\begin{note}[来源路径]

近世代数/近世代数/无标题/无标题 20bd2efa3f42803bb142e695e45bf369.md

\end{note}

\section*{第一题}

\begin{note}[来源上级主题]

3.9作业 (3\%209\%E4\%BD\%9C\%E4\%B8\%9A\%20205d2efa3f428077b8a0d0acdba61717.md)

\end{note}

\begin{homework}

证明1+i是极大理想 1+i可以表示成r（1+i）因为R是交换环且有单位元。

第一步证明1+i里面没有单位元

而a，b都是整数都是R里面的所以1在J里面所以J=R

c+di不在1+i理想里面 c+di-d(1+I)=c-d是整数且与1+i中2互素能组成一个1

剩余类

\end{homework}

\begin{note}[关联图片]
以下图片已纳入本条整理，当前版本优先依据 Markdown 文字整理，图片细节可在审核时对照：

1. 4601748744536\_.pic.jpg\\
2. image.png\\
3. image 1.png\\
4. image 2.png\\
5. image 3.png\\

\end{note}

\begin{note}[来源路径]

近世代数/近世代数/无标题/第一题 205d2efa3f4280d38576e58476d8c41d.md

\end{note}

\section*{第一题}

\begin{note}[来源上级主题]

3.2整环 (3\%202\%E6\%95\%B4\%E7\%8E\%AF\%20209d2efa3f4280cda320c843bba252f6.md)

\end{note}

\begin{homework}

(AeA Hl, oI RAR A ERLE

BF,

\end{homework}

\begin{note}[图片 OCR 摘录，待人工复核]
以下内容来自源图片识别，便于审核时对照：

\textbf{图片 1}（10961749123587\_.pic.jpg）
\begin{Verbatim}[fontsize=\small]
(AeA Hl, oI RAR A ERLE
BF,
\end{Verbatim}

\textbf{图片 2}（image.png）
\begin{Verbatim}[fontsize=\small]
Bor-y=0, WU:

_    ac + 2bd = 0
(ac + 2b) + (ad + boW/2 = 0 | ad + be =0
TBE OF i.
$8:V2¢Q BEM: pa +dv2 =08 BRE:
a=0, b=0
MATABRE:
-MRa=b=0, PRxe=0, $B
-MRc=d=0, Py=0, hee
Alt:
z,y#0>ayF#0
\end{Verbatim}

\end{note}

\begin{note}[来源路径]

近世代数/近世代数/无标题/第一题 209d2efa3f4280518e01dd4451714875.md

\end{note}

\section*{第三周作业}

\begin{note}[来源上级主题]

第二章作业 (\%E7\%AC\%AC\%E4\%BA\%8C\%E7\%AB\%A0\%E4\%BD\%9C\%E4\%B8\%9A\%2020bd2efa3f42805f9ecef7b781a02d72.md)

\end{note}

\begin{homework}

子级 项目: 第二题 (\%E7\%AC\%AC\%E4\%BA\%8C\%E9\%A2\%98\%2020bd2efa3f42809d8011e3ecfeda2cd1.md), 3，4 (3\%EF\%BC\%8C4\%2020bd2efa3f4280daa104ec7e9deedc65.md)

\end{homework}

\begin{note}[关联图片]
以下图片已纳入本条整理，当前版本优先依据 Markdown 文字整理，图片细节可在审核时对照：

1. image.png\\
2. image 1.png\\

\end{note}

\begin{note}[来源路径]

近世代数/近世代数/无标题/第三周作业 20bd2efa3f428003b74ac10b654adc8a.md

\end{note}

\section*{第三章作业}

\begin{homework}

子级 项目: 3.2整环 (3\%202\%E6\%95\%B4\%E7\%8E\%AF\%20209d2efa3f4280cda320c843bba252f6.md), 3.3 (3\%203\%20209d2efa3f428020b15ef2f886bfb11e.md), 3.4 (3\%204\%20209d2efa3f4280eca043fa61064594c9.md), 3.9作业 (3\%209\%E4\%BD\%9C\%E4\%B8\%9A\%20205d2efa3f428077b8a0d0acdba61717.md), 3.8作业 (3\%208\%E4\%BD\%9C\%E4\%B8\%9A\%20206d2efa3f42807aaea0cdba1c8e0355.md), 3.5作业 (3\%205\%E4\%BD\%9C\%E4\%B8\%9A\%20209d2efa3f4280eb969cc9dce9865299.md)

\end{homework}

\begin{note}[来源路径]

近世代数/近世代数/无标题/第三章作业 205d2efa3f428004b442cf8cf50564e6.md

\end{note}

\section*{第三题}

\begin{note}[来源上级主题]

3.9作业 (3\%209\%E4\%BD\%9C\%E4\%B8\%9A\%20205d2efa3f428077b8a0d0acdba61717.md)

\end{note}

\begin{homework}

3 RATE BTA SCIFI R, ELA) eR ARE AR) (W)Ae—- 1.

— 2 GI

ZIBEP |

\end{homework}

\begin{note}[图片 OCR 摘录，待人工复核]
以下内容来自源图片识别，便于审核时对照：

\textbf{图片 1}（9411748747842\_.pic.jpg）
\begin{Verbatim}[fontsize=\small]
3 RATE BTA SCIFI R, ELA) eR ARE AR) (W)Ae—- 1.
— 2 GI
ZIBEP |
\end{Verbatim}

\textbf{图片 2}（image.png）
\begin{Verbatim}[fontsize=\small]
M (4) HFS A—-W4in, Xn BBR. HUARR
Kji—P FHA, FFA C4) CU, (4) 4U. MA
WDSa= 2m#An
Hilba =4qt+ 2, MD4Aq+2-4q=2MA=(2)=R
RRB, (4) ER O-PRABA.
#2R/(4)#,02]4L0] fF £21] C23=C4)]=(0).
AI R/(4) BBR FARHAE—TPR,
\end{Verbatim}

\end{note}

\begin{note}[来源路径]

近世代数/近世代数/无标题/第三题 205d2efa3f4280c295f6ed5b008ce795.md

\end{note}

\section*{第三题}

\begin{note}[来源上级主题]

3.3 (3\%203\%20209d2efa3f428020b15ef2f886bfb11e.md)

\end{note}

\begin{homework}

除环就是找逆元

除环 需要单位元 大于一个元素 除了0有逆元

定理1除环只有两个理想0理想和单位理想

有限群每一个元素都有元

有限+无零因子=有逆元 有限可以让阶有限然后存在\$a\textasciicircum{}k\$=e

消去律成立与证明G是一个群 重点

\end{homework}

\begin{note}[关联图片]
以下图片已纳入本条整理，当前版本优先依据 Markdown 文字整理，图片细节可在审核时对照：

1. 11011749125713\_.pic.jpg\\
2. image.png\\
3. image 1.png\\

\end{note}

\begin{note}[来源路径]

近世代数/近世代数/无标题/第三题 209d2efa3f4280918485e1eeb5bf0217.md

\end{note}

\section*{第三题}

\begin{note}[来源上级主题]

第四章作业 (\%E7\%AC\%AC\%E5\%9B\%9B\%E7\%AB\%A0\%E4\%BD\%9C\%E4\%B8\%9A\%20204d2efa3f4280b4b282f638dd2dce07.md)

\end{note}

\begin{homework}

第一步先求单位 这里是正负1正负i

证明5不是不可约元 5有没有唯一分解？

5≠0 1/5不属于环

-1，1正负i都是单位，他的这些分解可以通过乘单位互相切换就是唯一分解

\end{homework}

\begin{note}[关联图片]
以下图片已纳入本条整理，当前版本优先依据 Markdown 文字整理，图片细节可在审核时对照：

1. image.png\\
2. image 1.png\\
3. image 2.png\\

\end{note}

\begin{note}[来源路径]

近世代数/近世代数/无标题/第三题- 204d2efa3f4280c58963d88dd6ba9b10.md

\end{note}

\section*{第三题再看}

\begin{note}[来源上级主题]

3.9作业 (3\%209\%E4\%BD\%9C\%E4\%B8\%9A\%20205d2efa3f428077b8a0d0acdba61717.md)

\end{note}

\begin{homework}

因为x是整数所以 c+d必须整除2所以 只能是[0]的剩余类 和1的剩余类

R/1+i的剩余类群 只有膜1+i=0和1两种情况

\end{homework}

\begin{note}[关联图片]
以下图片已纳入本条整理，当前版本优先依据 Markdown 文字整理，图片细节可在审核时对照：

1. 9391748746130\_.pic.jpg\\
2. image.png\\
3. image 1.png\\

\end{note}

\begin{note}[来源路径]

近世代数/近世代数/无标题/第三题再看 205d2efa3f42808f87efe6f9cc3fa0d4.md

\end{note}

\section*{第二章作业}

\begin{homework}

子级 项目: 第三周作业 (\%E7\%AC\%AC\%E4\%B8\%89\%E5\%91\%A8\%E4\%BD\%9C\%E4\%B8\%9A\%2020bd2efa3f428003b74ac10b654adc8a.md), 第四周作业 (\%E7\%AC\%AC\%E5\%9B\%9B\%E5\%91\%A8\%E4\%BD\%9C\%E4\%B8\%9A\%2020bd2efa3f4280a583d6e86ac317e9a1.md), 第五周作业 (\%E7\%AC\%AC\%E4\%BA\%94\%E5\%91\%A8\%E4\%BD\%9C\%E4\%B8\%9A\%2020bd2efa3f428007923bd76a797a9ede.md), 第六周子群 (\%E7\%AC\%AC\%E5\%85\%AD\%E5\%91\%A8\%E5\%AD\%90\%E7\%BE\%A4\%2020bd2efa3f4280cbaebef863c27b4e24.md), 找子群再看明天早上 (\%E6\%89\%BE\%E5\%AD\%90\%E7\%BE\%A4\%E5\%86\%8D\%E7\%9C\%8B\%E6\%98\%8E\%E5\%A4\%A9\%E6\%97\%A9\%E4\%B8\%8A\%2020bd2efa3f4280e39200d13ff510ae80.md), 循环群的子群还是循环群 (\%E5\%BE\%AA\%E7\%8E\%AF\%E7\%BE\%A4\%E7\%9A\%84\%E5\%AD\%90\%E7\%BE\%A4\%E8\%BF\%98\%E6\%98\%AF\%E5\%BE\%AA\%E7\%8E\%AF\%E7\%BE\%A4\%2020bd2efa3f42808aba03cfbc6170a459.md), 第五题找子群 (\%E7\%AC\%AC\%E4\%BA\%94\%E9\%A2\%98\%E6\%89\%BE\%E5\%AD\%90\%E7\%BE\%A4\%2020bd2efa3f4280cb98b6e5c0d60168df.md), 阶是素数的群一定是循环 (\%E9\%98\%B6\%E6\%98\%AF\%E7\%B4\%A0\%E6\%95\%B0\%E7\%9A\%84\%E7\%BE\%A4\%E4\%B8\%80\%E5\%AE\%9A\%E6\%98\%AF\%E5\%BE\%AA\%E7\%8E\%AF\%2020bd2efa3f428094bf7cca8d7e7f854f.md), 2正规子群 (2\%E6\%AD\%A3\%E8\%A7\%84\%E5\%AD\%90\%E7\%BE\%A4\%2020bd2efa3f4280879ad4e4e896b6e022.md), 陪集有正规子群自身 (\%E9\%99\%AA\%E9\%9B\%86\%E6\%9C\%89\%E6\%AD\%A3\%E8\%A7\%84\%E5\%AD\%90\%E7\%BE\%A4\%E8\%87\%AA\%E8\%BA\%AB\%2020bd2efa3f4280d5b2a5fb73f02da9c0.md), 循环群的商群还是循环群 (\%E5\%BE\%AA\%E7\%8E\%AF\%E7\%BE\%A4\%E7\%9A\%84\%E5\%95\%86\%E7\%BE\%A4\%E8\%BF\%98\%E6\%98\%AF\%E5\%BE\%AA\%E7\%8E\%AF\%E7\%BE\%A4\%2020bd2efa3f4280278080f5000566e0e2.md)

\end{homework}

\begin{note}[来源路径]

近世代数/近世代数/无标题/第二章作业 20bd2efa3f42805f9ecef7b781a02d72.md

\end{note}

\section*{第二题}

\begin{note}[来源上级主题]

3.9作业 (3\%209\%E4\%BD\%9C\%E4\%B8\%9A\%20205d2efa3f428077b8a0d0acdba61717.md)

\end{note}

\begin{homework}

BR j2 Yorso

CEA War:

Qbeod

61 Yeoro

—      —\textasciitilde{}Fh6

OED — Aha

v (sear    ae

f0d>= D

¥ o=0

\end{homework}

\begin{note}[图片 OCR 摘录，待人工复核]
以下内容来自源图片识别，便于审核时对照：

\textbf{图片 1}（9401748747439\_.pic.jpg）
\begin{Verbatim}[fontsize=\small]
BR j2 Yorso
CEA War:
Qbeod
61 Yeoro
—      —~Fh6
OED — Aha
v (sear    ae
f0d>= D
¥ o=0
\end{Verbatim}

\end{note}

\begin{note}[来源路径]

近世代数/近世代数/无标题/第二题 205d2efa3f4280d18289e2c61ebcba54.md

\end{note}

\section*{第二题}

\begin{note}[来源上级主题]

第四周作业 (\%E7\%AC\%AC\%E5\%9B\%9B\%E5\%91\%A8\%E4\%BD\%9C\%E4\%B8\%9A\%2020bd2efa3f4280a583d6e86ac317e9a1.md)

\end{note}

\begin{homework}

2 REABTARBERMRA. WE, WA ABT

RES

j      3-WH-6     0TIH7EEEEjo9E0

7E5EB9E3mXEEj3-XINEE2XiE4EFETE945EMEF7

n9GEE-WLEEEM5FEEEgEe-55EGE9Ef

Ts:         xvaxrt+b    O Hil b FEA ARM, a\#0

A;         zrcxt+d    c Ald Jk, c\#O

BA

tA;          xox? = (axrt+b)4= ¢ (axt+b) +d

= (ca)x+ (cb+d)

3X Hic a Fle b + d MELAMK, HHca+0.

PU cA HRTG.

EM RRB MIL, PU LR RM,

Fi AE HR

é;                   LL

JATG.

ES WMHE, «EG PAM, BD

wi,       xy x +(- 4)

A lLG 1 ER.

\{AG ARTA. \&

Tis                  zoxrt+l

Tos                   xL> 2x

BA

TiTes       a (e"!)"* = (@+1)"* =2n+2

\end{homework}

\begin{note}[图片 OCR 摘录，待人工复核]
以下内容来自源图片识别，便于审核时对照：

\textbf{图片 1}（image.png）
\begin{Verbatim}[fontsize=\small]
2 REABTARBERMRA. WE, WA ABT
RES
j      3-WH-6     0TIH7EEEEjo9E0
7E5EB9E3mXEEj3-XINEE2XiE4EFETE945EMEF7
n9GEE-WLEEEM5FEEEgEe-55EGE9Ef
Ts:         xvaxrt+b    O Hil b FEA ARM, a#0
A;         zrcxt+d    c Ald Jk, c#O
BA
tA;          xox? = (axrt+b)4= ¢ (axt+b) +d
= (ca)x+ (cb+d)
3X Hic a Fle b + d MELAMK, HHca+0.
PU cA HRTG.
EM RRB MIL, PU LR RM,
Fi AE HR
é;                   LL
JATG.
ES WMHE, «EG PAM, BD
wi,       xy x +(- 4)
A lLG 1 ER.
{AG ARTA. &
Tis                  zoxrt+l
Tos                   xL> 2x
BA
TiTes       a (e"!)"* = (@+1)"* =2n+2
\end{Verbatim}

\textbf{图片 2}（image 1.png）
\begin{Verbatim}[fontsize=\small]
TayT1     Ti Lor +1
TT it       x (x *)    (22x)
TT, FATT,
\end{Verbatim}

\end{note}

\begin{note}[来源路径]

近世代数/近世代数/无标题/第二题 20bd2efa3f42801e9db4fc42af376384.md

\end{note}

\section*{第二题}

\begin{note}[来源上级主题]

第三周作业 (\%E7\%AC\%AC\%E4\%B8\%89\%E5\%91\%A8\%E4\%BD\%9C\%E4\%B8\%9A\%2020bd2efa3f428003b74ac10b654adc8a.md)

\end{note}

\begin{homework}

\$a\textasciicircum{}n\$乘a逆的n次方-e 证明a逆的n次方也是e 要证明n就是a逆的阶然后再用一遍aa的逆=e然后证明a，a的逆和任意不和他们相等的b产生的b的逆也不相等

看见阶是2一定要想到自己等于自己的逆

\end{homework}

\begin{note}[关联图片]
以下图片已纳入本条整理，当前版本优先依据 Markdown 文字整理，图片细节可在审核时对照：

1. image.png\\

\end{note}

\begin{note}[来源路径]

近世代数/近世代数/无标题/第二题 20bd2efa3f42809d8011e3ecfeda2cd1.md

\end{note}

\section*{第五周作业}

\begin{note}[来源上级主题]

第二章作业 (\%E7\%AC\%AC\%E4\%BA\%8C\%E7\%AB\%A0\%E4\%BD\%9C\%E4\%B8\%9A\%2020bd2efa3f42805f9ecef7b781a02d72.md)

\end{note}

\begin{homework}

子级 项目: 循环群 (\%E5\%BE\%AA\%E7\%8E\%AF\%E7\%BE\%A4\%2020bd2efa3f4280d99b3eea67ac6b65ad.md), 4 (4\%2020bd2efa3f4280efba7cd41e90a2076b.md)

\end{homework}

\begin{note}[来源路径]

近世代数/近世代数/无标题/第五周作业 20bd2efa3f428007923bd76a797a9ede.md

\end{note}

\section*{第五题找子群}

\begin{note}[来源上级主题]

第二章作业 (\%E7\%AC\%AC\%E4\%BA\%8C\%E7\%AB\%A0\%E4\%BD\%9C\%E4\%B8\%9A\%2020bd2efa3f42805f9ecef7b781a02d72.md)

\end{note}

\begin{homework}

5. FRM MAA MDT A.

RB Pi2HMRAMHCR-THALMEMR. Alt

Hi 4, GHFRMRLGAH. ABB:

(C01) =£0]

(CiDy=(CS5)=C7)= (CUD) =E

(C27) = (€10]) =\{£ 2], (41, [6], [8], [10], COI\}

((3])=((9]) =\{[31, [6], £91, COI\}

(4) =(C8)=\{041, £81, COT\}

([6])= \{£6], £07\}

EGH HATH.

\end{homework}

\begin{note}[图片 OCR 摘录，待人工复核]
以下内容来自源图片识别，便于审核时对照：

\textbf{图片 1}（image.png）
\begin{Verbatim}[fontsize=\small]
5. FRM MAA MDT A.

RB Pi2HMRAMHCR-THALMEMR. Alt
Hi 4, GHFRMRLGAH. ABB:

(C01) =£0]

(CiDy=(CS5)=C7)= (CUD) =E

(C27) = (€10]) ={£ 2], (41, [6], [8], [10], COI}

((3])=((9]) ={[31, [6], £91, COI}

(4) =(C8)={041, £81, COT}

([6])= {£6], £07}
EGH HATH.
\end{Verbatim}

\end{note}

\begin{note}[来源路径]

近世代数/近世代数/无标题/第五题找子群 20bd2efa3f4280cb98b6e5c0d60168df.md

\end{note}

\section*{第六周子群}

\begin{note}[来源上级主题]

第二章作业 (\%E7\%AC\%AC\%E4\%BA\%8C\%E7\%AB\%A0\%E4\%BD\%9C\%E4\%B8\%9A\%2020bd2efa3f42805f9ecef7b781a02d72.md)

\end{note}

\begin{homework}

待根据图片进一步人工校对。

\end{homework}

\begin{note}[来源路径]

近世代数/近世代数/无标题/第六周子群 20bd2efa3f4280cbaebef863c27b4e24.md

\end{note}

\section*{第四周作业}

\begin{note}[来源上级主题]

第二章作业 (\%E7\%AC\%AC\%E4\%BA\%8C\%E7\%AB\%A0\%E4\%BD\%9C\%E4\%B8\%9A\%2020bd2efa3f42805f9ecef7b781a02d72.md)

\end{note}

\begin{homework}

子级 项目: 第二题 (\%E7\%AC\%AC\%E4\%BA\%8C\%E9\%A2\%98\%2020bd2efa3f42801e9db4fc42af376384.md), 置换群 (\%E7\%BD\%AE\%E6\%8D\%A2\%E7\%BE\%A4\%2020bd2efa3f42808c9256d550e77037b7.md), 无标题 (\%E6\%97\%A0\%E6\%A0\%87\%E9\%A2\%98\%2020bd2efa3f42803bb142e695e45bf369.md)

\end{homework}

\begin{note}[来源路径]

近世代数/近世代数/无标题/第四周作业 20bd2efa3f4280a583d6e86ac317e9a1.md

\end{note}

\section*{第四章作业}

\begin{homework}

子级 项目: 第三题- (\%E7\%AC\%AC\%E4\%B8\%89\%E9\%A2\%98-\%20204d2efa3f4280c58963d88dd6ba9b10.md), 1欧式环 (1\%E6\%AC\%A7\%E5\%BC\%8F\%E7\%8E\%AF\%20204d2efa3f42806fb730dc64fc6e2b23.md)

\end{homework}

\begin{note}[来源路径]

近世代数/近世代数/无标题/第四章作业 204d2efa3f4280b4b282f638dd2dce07.md

\end{note}

\section*{置换群}

\begin{note}[来源上级主题]

第四周作业 (\%E7\%AC\%AC\%E5\%9B\%9B\%E5\%91\%A8\%E4\%BD\%9C\%E4\%B8\%9A\%2020bd2efa3f4280a583d6e86ac317e9a1.md)

\end{note}

\begin{homework}

置换是一一对应的变化，不相连就是交集是空，写成循环置换就是为了省地方是按照特定顺序的

n次对称群就是全体置换

置换是一个动作 是特殊的变换为了不写这么长才发明出循环置换这里就135循环 24循环写成(1 3 5)*(2 4)

有限群都与置换群同构

证明是逆还是要证明乘起来是单位

\end{homework}

\begin{note}[关联图片]
以下图片已纳入本条整理，当前版本优先依据 Markdown 文字整理，图片细节可在审核时对照：

1. image.png\\
2. image 1.png\\
3. image 2.png\\
4. image 3.png\\

\end{note}

\begin{note}[来源路径]

近世代数/近世代数/无标题/置换群 20bd2efa3f42808c9256d550e77037b7.md

\end{note}

\section*{阶是素数的群一定是循环}

\begin{note}[来源上级主题]

第二章作业 (\%E7\%AC\%AC\%E4\%BA\%8C\%E7\%AB\%A0\%E4\%BD\%9C\%E4\%B8\%9A\%2020bd2efa3f42805f9ecef7b781a02d72.md)

\end{note}

\begin{homework}

也就是阶是素数的群每个不是单位元的都是生成元 先生成一个循环子群子群的阶是父群阶的因子 因为父群阶是素数所以子群阶就是父群的阶

\end{homework}

\begin{note}[关联图片]
以下图片已纳入本条整理，当前版本优先依据 Markdown 文字整理，图片细节可在审核时对照：

1. image.png\\

\end{note}

\begin{note}[来源路径]

近世代数/近世代数/无标题/阶是素数的群一定是循环 20bd2efa3f428094bf7cca8d7e7f854f.md

\end{note}

\section*{陪集有正规子群自身}

\begin{note}[来源上级主题]

第二章作业 (\%E7\%AC\%AC\%E4\%BA\%8C\%E7\%AB\%A0\%E4\%BD\%9C\%E4\%B8\%9A\%2020bd2efa3f42805f9ecef7b781a02d72.md)

\end{note}

\begin{homework}

3. TER, AR 2H TR EAARTH.

SGE—-+RIN EGH—-TBRA 2TH.

HneEN, PABRAnN=Nn. KHOEG, SEN.

PATNA 2, GRAMBITAREN MON; G

ERED RPT ASN ALN 6. AON =Nb, XR WF

Gites oR, @GN=aNofNEGH—-+ABTER.,

\end{homework}

\begin{note}[图片 OCR 摘录，待人工复核]
以下内容来自源图片识别，便于审核时对照：

\textbf{图片 1}（image.png）
\begin{Verbatim}[fontsize=\small]
3. TER, AR 2H TR EAARTH.

 SGE—-+RIN EGH—-TBRA 2TH.

HneEN, PABRAnN=Nn. KHOEG, SEN.
PATNA 2, GRAMBITAREN MON; G
ERED RPT ASN ALN 6. AON =Nb, XR WF
Gites oR, @GN=aNofNEGH—-+ABTER.,
\end{Verbatim}

\end{note}

\begin{note}[来源路径]

近世代数/近世代数/无标题/陪集有正规子群自身 20bd2efa3f4280d5b2a5fb73f02da9c0.md

\end{note}
